\subsection{List Mutation Operations}

A list is mutable. This means that the list object can be changed in place after it has been created. New references can be added, existing slots can be replaced, ranges can be overwritten, elements can be removed, and the current order can be changed.

\begin{verbatim}
items = ["spam", "eggs"]

print(f"before: {items!r}")
print(f"id:     {id(items)}")

items.append(42)

print()
print(f"after:  {items!r}")
print(f"id:     {id(items)}")
\end{verbatim}

Possible output:

\begin{verbatim}
before: ['spam', 'eggs']
id:     140619786870656

after:  ['spam', 'eggs', 42]
id:     140619786870656
\end{verbatim}

The identity of the list is the same. The object was mutated.

This is different from string operations, where a new string object is created because strings are immutable.

\begin{verbatim}
text = "spam"
items = ["spam"]

text_id = id(text)
items_id = id(items)

text = text + " eggs"
items.append("eggs")

print(f"text:  {text!r}")
print(f"items: {items!r}")

print()
print(f"text same object:  {id(text) == text_id}")
print(f"items same object: {id(items) == items_id}")
\end{verbatim}

Possible output:

\begin{verbatim}
text:  'spam eggs'
items: ['spam', 'eggs']

text same object:  False
items same object: True
\end{verbatim}

The methods in this section are the ordinary mutation tools used with lists.

\begin{verbatim}
items = ["Homer", "Marge", "Bart"]

selected_names = [
    "append",
    "extend",
    "insert",
    "remove",
    "pop",
    "sort",
    "reverse",
    "copy",
]

print(f"type(items): {type(items)}")
print()

for name in selected_names:
    print(f"{name:<10} {name in dir(items)}")
\end{verbatim}

Possible output:

\begin{verbatim}
type(items): <class 'list'>

append     True
extend     True
insert     True
remove     True
pop        True
sort       True
reverse    True
copy       True
\end{verbatim}

\subsubsection{Appending and Extending: \texttt{.append()} vs. \texttt{.extend()}}

The method \texttt{.append()} adds one object to the end of the list.

\begin{verbatim}
items = ["spam", "eggs"]

items.append("bacon")
items.append(42)

print(items)
\end{verbatim}

Possible output:

\begin{verbatim}
['spam', 'eggs', 'bacon', 42]
\end{verbatim}

The argument passed to \texttt{append()} is inserted as one element. If the argument is itself a list, that whole list becomes one element.

\begin{verbatim}
items = ["spam", "eggs"]
more = ["bacon", "beans"]

items.append(more)

print(f"items: {items!r}")
print(f"len(items): {len(items)}")

print()
for index, item in enumerate(items):
    print(f"index={index} value={item!r} type={type(item)}")
\end{verbatim}

Possible output:

\begin{verbatim}
items: ['spam', 'eggs', ['bacon', 'beans']]
len(items): 3

index=0 value='spam' type=<class 'str'>
index=1 value='eggs' type=<class 'str'>
index=2 value=['bacon', 'beans'] type=<class 'list'>
\end{verbatim}

The list \texttt{more} was not unpacked. It was appended as one object.

The method \texttt{.extend()} takes an iterable and appends its elements one by one.

\begin{verbatim}
items = ["spam", "eggs"]
more = ["bacon", "beans"]

items.extend(more)

print(f"items: {items!r}")
print(f"len(items): {len(items)}")

print()
for index, item in enumerate(items):
    print(f"index={index} value={item!r} type={type(item)}")
\end{verbatim}

Possible output:

\begin{verbatim}
items: ['spam', 'eggs', 'bacon', 'beans']
len(items): 4

index=0 value='spam' type=<class 'str'>
index=1 value='eggs' type=<class 'str'>
index=2 value='bacon' type=<class 'str'>
index=3 value='beans' type=<class 'str'>
\end{verbatim}

So the practical distinction is:

\begin{itemize}
    \item \texttt{append(x)} adds \texttt{x} as one new element;
    \item \texttt{extend(xs)} walks through \texttt{xs} and adds each component.
\end{itemize}

This difference is especially visible with strings.

\begin{verbatim}
items1 = ["Ni"]
items2 = ["Ni"]

items1.append("abc")
items2.extend("abc")

print(f"items1: {items1!r}")
print(f"items2: {items2!r}")
\end{verbatim}

Possible output:

\begin{verbatim}
items1: ['Ni', 'abc']
items2: ['Ni', 'a', 'b', 'c']
\end{verbatim}

A string is iterable, so \texttt{extend("abc")} appends its characters one by one. Since Python has no separate character type, each appended character is a one-character \texttt{str} object.

Both methods mutate the same list object.

\begin{verbatim}
items = ["Jerry"]

old_id = id(items)

items.append("Elaine")
items.extend(["George", "Kramer"])

print(f"items: {items!r}")
print(f"same object: {id(items) == old_id}")
\end{verbatim}

Possible output:

\begin{verbatim}
items: ['Jerry', 'Elaine', 'George', 'Kramer']
same object: True
\end{verbatim}

These methods return \texttt{None}. The result is stored in the mutated list, not in the method return value.

\begin{verbatim}
items = ["Homer"]

result = items.append("Marge")

print(f"items:  {items!r}")
print(f"result: {result!r}")
\end{verbatim}

Possible output:

\begin{verbatim}
items:  ['Homer', 'Marge']
result: None
\end{verbatim}

A common mistake is therefore:

\begin{verbatim}
items = ["Homer"]
items = items.append("Marge")

print(items)
\end{verbatim}

Possible output:

\begin{verbatim}
None
\end{verbatim}

The name \texttt{items} was rebound to the return value of \texttt{append()}, which is \texttt{None}. The correct form is to call \texttt{append()} without assigning its return value.

\subsubsection{Index Assignment and Slice Assignment}

Index assignment replaces the object reference stored at one list position.

\begin{verbatim}
items = ["spam", "eggs", "bacon"]

print(f"before: {items!r}")

items[1] = "beans"

print(f"after:  {items!r}")
\end{verbatim}

Possible output:

\begin{verbatim}
before: ['spam', 'eggs', 'bacon']
after:  ['spam', 'beans', 'bacon']
\end{verbatim}

The list object stays the same. One slot now points to another object.

\begin{verbatim}
items = ["spam", "eggs", "bacon"]

old_id = id(items)
old_element = items[1]

items[1] = "beans"

print(f"items: {items!r}")
print(f"same list object: {id(items) == old_id}")

print()
print(f"old_element: {old_element!r}")
print(f"new element: {items[1]!r}")
\end{verbatim}

Possible output:

\begin{verbatim}
items: ['spam', 'beans', 'bacon']
same list object: True

old_element: 'eggs'
new element: 'beans'
\end{verbatim}

An invalid index assignment raises \texttt{IndexError}.

\begin{verbatim}
items = ["spam", "eggs"]

try:
    items[5] = "bacon"
except IndexError as err:
    print("assignment failed")
    print(f"message: {err}")
\end{verbatim}

Possible output:

\begin{verbatim}
assignment failed
message: list assignment index out of range
\end{verbatim}

Index assignment replaces an existing position. It does not create missing positions.

Slice assignment replaces a range of positions.

\begin{verbatim}
items = ["a", "b", "c", "d", "e"]

print(f"before: {items!r}")

items[1:4] = ["X", "Y"]

print(f"after:  {items!r}")
\end{verbatim}

Possible output:

\begin{verbatim}
before: ['a', 'b', 'c', 'd', 'e']
after:  ['a', 'X', 'Y', 'e']
\end{verbatim}

The slice \texttt{items[1:4]} selected \texttt{"b"}, \texttt{"c"}, and \texttt{"d"}. Those three elements were replaced by two new elements.

Slice assignment can therefore change the length of the list.

\begin{verbatim}
items = ["a", "b", "c"]

print(f"before: {items!r}")
print(f"len before: {len(items)}")

items[1:2] = ["X", "Y", "Z"]

print()
print(f"after: {items!r}")
print(f"len after: {len(items)}")
\end{verbatim}

Possible output:

\begin{verbatim}
before: ['a', 'b', 'c']
len before: 3

after: ['a', 'X', 'Y', 'Z', 'c']
len after: 5
\end{verbatim}

An empty slice can be used as an insertion point.

\begin{verbatim}
items = ["a", "d"]

items[1:1] = ["b", "c"]

print(items)
\end{verbatim}

Possible output:

\begin{verbatim}
['a', 'b', 'c', 'd']
\end{verbatim}

The slice \texttt{items[1:1]} selects no existing elements, but it identifies the boundary before index \texttt{1}.

Slice assignment requires an iterable on the right side.

\begin{verbatim}
items = ["a", "b", "c"]

try:
    items[1:2] = 42
except TypeError as err:
    print("slice assignment failed")
    print(f"message: {err}")
\end{verbatim}

Possible output:

\begin{verbatim}
slice assignment failed
message: can only assign an iterable
\end{verbatim}

If a string is assigned to a slice, the string is iterated character by character.

\begin{verbatim}
items = ["a", "b", "c"]

items[1:2] = "XY"

print(items)
\end{verbatim}

Possible output:

\begin{verbatim}
['a', 'X', 'Y', 'c']
\end{verbatim}

This is the same iterable rule seen with \texttt{extend()}.

The practical rule is:

\begin{quote}
Index assignment replaces one existing slot. Slice assignment replaces a range of slots and may also change the list length.
\end{quote}

\subsubsection{Removing Elements: \texttt{.remove()}, \texttt{.pop()}, and \texttt{del}}

Lists provide several removal mechanisms. They differ by how the target is chosen.

\begin{itemize}
    \item \texttt{remove(value)} removes the first equal value;
    \item \texttt{pop(index)} removes by position and returns the removed object;
    \item \texttt{del items[index]} removes by position without returning the object;
    \item \texttt{del items[start:stop]} removes a slice.
\end{itemize}

The method \texttt{.remove()} searches by value.

\begin{verbatim}
items = ["spam", "eggs", "spam", "bacon"]

items.remove("spam")

print(items)
\end{verbatim}

Possible output:

\begin{verbatim}
['eggs', 'spam', 'bacon']
\end{verbatim}

Only the first matching value was removed.

If the value is absent, \texttt{remove()} raises \texttt{ValueError}.

\begin{verbatim}
items = ["spam", "eggs"]

try:
    items.remove("bacon")
except ValueError as err:
    print("remove failed")
    print(f"message: {err}")
\end{verbatim}

Possible output:

\begin{verbatim}
remove failed
message: list.remove(x): x not in list
\end{verbatim}

The method \texttt{.pop()} removes by index and returns the removed object. With no argument, it removes the last element.

\begin{verbatim}
items = ["Jerry", "Elaine", "George", "Kramer"]

last = items.pop()

print(f"last:  {last!r}")
print(f"items: {items!r}")
\end{verbatim}

Possible output:

\begin{verbatim}
last:  'Kramer'
items: ['Jerry', 'Elaine', 'George']
\end{verbatim}

With an explicit index, \texttt{pop()} removes that position.

\begin{verbatim}
items = ["Jerry", "Elaine", "George", "Kramer"]

removed = items.pop(1)

print(f"removed: {removed!r}")
print(f"items:   {items!r}")
\end{verbatim}

Possible output:

\begin{verbatim}
removed: 'Elaine'
items:   ['Jerry', 'George', 'Kramer']
\end{verbatim}

An invalid pop index raises \texttt{IndexError}.

\begin{verbatim}
items = ["Jerry"]

try:
    items.pop(42)
except IndexError as err:
    print("pop failed")
    print(f"message: {err}")
\end{verbatim}

Possible output:

\begin{verbatim}
pop failed
message: pop index out of range
\end{verbatim}

The statement \texttt{del} removes a list element by position.

\begin{verbatim}
items = ["Homer", "Marge", "Bart", "Lisa"]

del items[2]

print(items)
\end{verbatim}

Possible output:

\begin{verbatim}
['Homer', 'Marge', 'Lisa']
\end{verbatim}

Unlike \texttt{pop()}, \texttt{del} does not return the removed object. It is a statement, not a method call.

\begin{verbatim}
items = ["Homer", "Marge", "Bart", "Lisa", "Maggie"]

del items[1:4]

print(items)
\end{verbatim}

Possible output:

\begin{verbatim}
['Homer', 'Maggie']
\end{verbatim}

The slice \texttt{items[1:4]} selected \texttt{"Marge"}, \texttt{"Bart"}, and \texttt{"Lisa"}, and \texttt{del} removed that range.

The three removal styles are often chosen like this:

\begin{verbatim}
items = ["a", "b", "c", "d"]

# remove by known value
items.remove("b")

# remove by position and keep removed object
removed = items.pop(1)

# remove by position and ignore removed object
del items[0]

print(f"removed: {removed!r}")
print(f"items:   {items!r}")
\end{verbatim}

Possible output:

\begin{verbatim}
removed: 'c'
items:   ['d']
\end{verbatim}

The practical rule is:

\begin{quote}
Use \texttt{remove()} when the value identifies the target. Use \texttt{pop()} when the position identifies the target and the removed object is needed. Use \texttt{del} when the position or slice identifies the target and no returned object is needed.
\end{quote}

\subsubsection{Sorting and Reversing In Place: \texttt{.sort()} and \texttt{.reverse()}}

The method \texttt{.sort()} sorts a list in place.

\begin{verbatim}
items = [30, 10, 42, 20]

old_id = id(items)

result = items.sort()

print(f"items:  {items!r}")
print(f"result: {result!r}")
print(f"same object: {id(items) == old_id}")
\end{verbatim}

Possible output:

\begin{verbatim}
items:  [10, 20, 30, 42]
result: None
same object: True
\end{verbatim}

The method changes the list and returns \texttt{None}. This follows the same mutation convention as \texttt{append()} and \texttt{extend()}.

The method \texttt{.reverse()} also mutates the list in place.

\begin{verbatim}
items = ["Jerry", "Elaine", "George", "Kramer"]

old_id = id(items)

result = items.reverse()

print(f"items:  {items!r}")
print(f"result: {result!r}")
print(f"same object: {id(items) == old_id}")
\end{verbatim}

Possible output:

\begin{verbatim}
items:  ['Kramer', 'George', 'Elaine', 'Jerry']
result: None
same object: True
\end{verbatim}

Sorting requires the elements to be comparable.

\begin{verbatim}
items = [10, "twenty", 30]

try:
    items.sort()
except TypeError as err:
    print("sort failed")
    print(f"message: {err}")
\end{verbatim}

Possible output:

\begin{verbatim}
sort failed
message: '<' not supported between instances of 'str' and 'int'
\end{verbatim}

The list may hold heterogeneous objects, but sorting still needs a meaningful ordering between the compared elements.

A \texttt{key} function can tell Python what value to sort by. For simple cases, built-in functions are enough.

\begin{verbatim}
names = ["Kramer", "Ed", "George", "Elaine", "Bob"]

names.sort(key=len)

print(names)
\end{verbatim}

Possible output:

\begin{verbatim}
['Ed', 'Bob', 'Kramer', 'George', 'Elaine']
\end{verbatim}

The list was sorted by string length.

Sorting can also be reversed.

\begin{verbatim}
numbers = [10, 30, 20, 42]

numbers.sort(reverse=True)

print(numbers)
\end{verbatim}

Possible output:

\begin{verbatim}
[42, 30, 20, 10]
\end{verbatim}

For strings, the default order is lexicographic and case-sensitive.

\begin{verbatim}
words = ["spam", "Eggs", "bacon", "Beans"]

words.sort()

print(words)
\end{verbatim}

Possible output:

\begin{verbatim}
['Beans', 'Eggs', 'bacon', 'spam']
\end{verbatim}

A case-normalizing key can make the ordering less surprising for ordinary text.

\begin{verbatim}
words = ["spam", "Eggs", "bacon", "Beans"]

words.sort(key=str.casefold)

print(words)
\end{verbatim}

Possible output:

\begin{verbatim}
['bacon', 'Beans', 'Eggs', 'spam']
\end{verbatim}

The function \texttt{sorted()} is different from \texttt{.sort()}. It creates a new sorted list and leaves the original object unchanged.

\begin{verbatim}
items = [30, 10, 42, 20]

new_items = sorted(items)

print(f"items:     {items!r}")
print(f"new_items: {new_items!r}")

print()
print(f"same object: {items is new_items}")
\end{verbatim}

Possible output:

\begin{verbatim}
items:     [30, 10, 42, 20]
new_items: [10, 20, 30, 42]

same object: False
\end{verbatim}

The same distinction exists between \texttt{.reverse()} and \texttt{reversed()}.

\begin{verbatim}
items = ["a", "b", "c"]

new_items = list(reversed(items))

print(f"items:     {items!r}")
print(f"new_items: {new_items!r}")
\end{verbatim}

Possible output:

\begin{verbatim}
items:     ['a', 'b', 'c']
new_items: ['c', 'b', 'a']
\end{verbatim}

The practical rule is:

\begin{quote}
\texttt{.sort()} and \texttt{.reverse()} mutate the existing list and return \texttt{None}. Use \texttt{sorted()} or \texttt{reversed()} when a separate result is wanted.
\end{quote}